\SourceLeaf{231}%
I meant here only to give the outlines of the proof, without any of the refinement or polish that I have no time to give it now. You are free to communicate it to your son. In any case, I ask you to greet him warmly for me and assure him of my particular esteem. At the same time, though, urge him to take up the following problem.

»To determine the volume of the tetrahedron (the space bounded by four planes)«.

Since the area of a triangle can be given so simply, one would have expected an equally simple expression for this volume. Yet this expectation, it seems, is disappointed.

To treat geometry properly from the beginning, it is indispensable to prove the possibility of a plane. The usual definition contains too much and in fact already implies a theorem surreptitiously. It is astonishing that all writers from Euclid down to the present have proceeded so carelessly in this matter. Yet this difficulty is of an altogether different [nature] from that of deciding between $\Sigma$ and $S$, and the former is not all that difficult to resolve. Your book will probably already satisfy me on this point as well.

It is precisely the impossibility of deciding between $\Sigma$ and $S$ a priori that gives the clearest proof that Kant was wrong to claim that space is only the \textls{form} of our intuition. I indicated another equally strong reason in a short essay that appeared in the Göttingischen Gelehrten Anzeigen in 1831, issue 64, pag.~625. You may not regret taking the trouble to obtain this volume of the G.G.A. (which any bookseller in Vienna or Ofen can easily arrange), since among other things it also sets out the quintessence of my view of imaginary quantities in a couple of pages. \mbox{.\ .\ .\ .\ .\ .\ .}

